\scanpage{105}
\mspage{92}
% archive sheet number: 104

\uline{Démonstration} Par 4.6, on sait déjà que
$(M',W',C')$ est une cat. à cofibrations spéciale.
[\uline{NB} Ce qui est noté $C'$ ici, était noté $\widetilde{C'}$ dans
la prop.~4.5. On doit utiliser la dernière
assertion de 4.6, où on suppose $C$ stable par
facteurs directs, OK car $(W,C,F)$ supposé
clos] Pour prouver que $(W',C',F'\overset{\mathrm{déf}}=(TC')_{*})$
est un triple de Quillen spécial, il
reste, par 5.3. et définition 5.6, à prouver
qu'il existe une petite partie $TC'_{0}$ de
$TC'_{\unsure{9|0}}$, telle que
\[
 (TC'_{0}\unsure{\subset|\subseteq})\quad TC'\subset (TC'_{0})^{\sim}. \tag*{(5.8.6)}
\]
(Le triple de Quillen sera de plus clos, puisque
$\deleted{\unsure{T}}C'$ est \inserted{\unsure{Q-}}clos \inserted{\unsure{à g.|à g}} par définition.) Soit donc
$TC_{0}$ une petite partie de $TC$ telle que
\[
 TC\subset (TC_{0})^{\sim},\tag*{(5.8.7)}
\]
qui existe par hypothèse sur $(W,C,F)$. On
prendra
\[
 TC'_{0}\overset{\mathrm{déf}}=u_{!}(TC_{0}).\tag*{(5.8.8)}
\]
On a vu dans 4.7. qu'on a
\[
 u_{!}(TC)\subset TC'\tag*{(5.8.9)}
\]
(\unsure{résulte} de l'hyp. $u^{*}(u_{!}(TC))\subset TC$
de (5.8.2 c)), \unsure{a|à} fortiori
\[
 TC'_{0}=u_{!}(TC_{0})\subset TC',
\]

